\section{Qualification, representation, and attained economic cost}\label{sec:study81}
The preceding attained-controller experiment used a coefficientwise qualification bound. Its ReLU--Bernstein ordering is an observation about those specified constructions. We now examine whether that ordering persists after the two bases receive the same qualification charge on their common affine functions. This directly tests a source of representation-dependent conservatism while preserving the original nonlinear investment economy and the Neural Bellman Operator construction.

\subsection{A prospective factorial comparison}
The design fixes two bases, two qualification rules and two acquisition resolutions. The eight constructed controllers and the installed zero controller are evaluated on the same original-law paths. Each basis has nine coefficients with absolute bounds eight, the same node covers, thirty-three-action search menus and downward action quantum $1/4096$. The calibrated rule replaces the coefficientwise modulus by \eqref{eq:relu-modulus81} or \eqref{eq:bernstein-modulus81}, both in residual qualification and in the regularity allowance passed to the previous fitted date. Both own-future continuations are constructed backwards. An exact primal--dual gap of at most $1/65536$ qualifies each new fit. The coefficient proposal is quantized at forty bits; no hidden-feature optimizer or outcome-selected basis is introduced.

A source manifest was committed before construction, and the complete actor identities were committed before generating 524,288 new four-coordinate continuous-bin rows. The bin precision is forty bits. Two estimators, cash cost and exact last-innovation conditional expectation, evaluate the same expected costs under the original uniform initial law and continuous innovations. For each estimator the frozen family contains nine marginal costs, thirty-six paired differences and two basis-by-qualification interactions. All ninety-four estimands share a single failure allowance bounded by $376\exp(-11)<1/100$. Shared paths and the two estimators do not multiply the number of independent rows.

The whole-bin enclosures are independently computed in outward binary64 arithmetic and directed-integer arithmetic. Their hull, intersected with the proved support, is used for inference. An exact rational implementation additionally checks 512 midpoint trajectories per controller. All nine controllers, bins, endpoints, exact sample moments and disagreements are part of the retained protocol. This revision re-evaluates every bin and every inference quantity; it neither regenerates those observations as a second sample nor substitutes replay time for the original service clock.

\begin{table}[htbp]
\centering
\caption{Basis comparison before and after qualification calibration}\label{tab:basis81}
{\small\setlength{\tabcolsep}{3.5pt}
\begin{tabular}{lrcc}
\toprule
Qualification & $n$ & ReLU minus Bernstein & Lower cost \\
\midrule
Legacy & 8 & [-0.077633, -0.075733] & R \\
Legacy & 16 & [-0.077232, -0.075307] & R \\
Calibrated & 8 & [0.003571, 0.004688] & B \\
Calibrated & 16 & [0.000955, 0.002013] & B \\
\bottomrule\end{tabular}}
\par\smallskip{\footnotesize All entries are simultaneous paired original-law cost intervals from the same prospective R80 family, with exact last-innovation conditional integration. Positive values favor Bernstein. Equal coefficient counts and affine-exact modulus charging do not make the two nonlinear function classes identical. The old fits are evaluated, not overwritten.}
\end{table}

\begin{table}[htbp]
\centering
\caption{Actual controller gains from calibrated qualification}\label{tab:calibration-gains81}
{\small\setlength{\tabcolsep}{3.5pt}
\begin{tabular}{lrc}
\toprule
Basis & $n$ & Legacy minus calibrated cost \\
\midrule
Relu & 8 & [0.000457, 0.001653] \\
Relu & 16 & [0.005159, 0.006341] \\
Bernstein & 8 & [0.080964, 0.082772] \\
Bernstein & 16 & [0.082601, 0.084405] \\
\bottomrule\end{tabular}}
\par\smallskip{\footnotesize Each comparison retains its basis, coefficient count and bounds, robust action menus, observation resolution and original economic primitives. Both attained future fits are reconstructed under the stated qualification rule. The displayed gains are not timing estimates or claims about unrestricted neural fitting.}
\end{table}


The calibration improves the expected cost of each basis at each resolution. In particular, the fine ReLU controller gains between $0.005159$ and $0.006341$ relative to its legacy construction. The matched basis ordering nevertheless reverses: calibrated Bernstein has lower cost than calibrated ReLU at both resolutions. At $n=16$ the ReLU-minus-Bernstein interval is $[0.000955,0.002013]$, rounded outward. That strict ordering coexists with equivalence at the declared margin $1/256$; a small positive difference can satisfy both statements. The corresponding coarse-resolution interval does not establish equivalence at that margin.

\begin{table}[htbp]
\centering
\caption{Predeclared basis-by-qualification interactions}\label{tab:interactions81}
{\small\setlength{\tabcolsep}{3.5pt}
\begin{tabular}{rc}
\toprule
$n$ & Change in Bernstein-minus-ReLU cost \\
\midrule
8 & [-0.082254, -0.079371] \\
16 & [-0.079209, -0.076297] \\
\bottomrule\end{tabular}}
\par\smallskip{\footnotesize The interaction is (calibrated Bernstein minus calibrated ReLU) minus (legacy Bernstein minus legacy ReLU). These four-term contrasts were included in the source-frozen family before the new cost rows were generated. They are not obtained by treating the four marginal estimates as independent.}
\end{table}


The predeclared interactions isolate the change in the basis comparison. They do not establish that the bases span identical functions, that the Bernstein derivative-control bound is globally exact, or that every other numerical difference has been removed. The positive operational conclusion is more specific: charging the function's derivative structure changes the attained original-law policy cost, and a policy ranking cannot be attributed to a representation name without examining its qualification rule. The historical ReLU gain and the new reversed matched ranking remain separate valid observations.

\subsection{Precision of the economic comparison}
The target was a paired total width of $1/256$, fixed before seeing the new rows. Appendix~\ref{app:factorial81} reports every interval and its width components, including every marginal interval whether or not its width falls below that paired target. Variance reduction by conditional integration affects the variance component, not the bounded-support penalty or the acquisition contract. For example, in the fine calibrated basis comparison the conditional interval's width consists of approximately $0.001008605$ from the finite-support term, $0.000048152$ from endpoint sampling-variance allowances, and $2.69\times10^{-11}$ from the numerical enclosure. The support term is therefore the principal remaining contribution in that comparison.

The numerical component includes the implemented handling of continuous-bin width, acquisition ambiguity and arithmetic propagation. It is not represented as a pure floating-point error estimate. Likewise, a small numerical component does not justify dropping the finite-sample support term. These diagnostics explain the achieved precision without inferring significance from raw mean differences or treating the inherited R67 policy-cost rows as observations of these newly constructed actors.

\subsection{From a fitted readout to a qualified optimizer class}
For each new terminal-label qualification problem the action-switch catalogue covers every first-date observation cell and all thirty-two alternative robust actions. The 20,480 rational dual bounds are independently recomputed using the antiderivative and box-infimum calculations of Appendix~\ref{app:dual81}; no switch optimization is rerun to replace an unfavorable certificate.

\begin{table}[htbp]
\centering
\caption{First-date action stability of qualification sublevels}\label{tab:stability81}
{\small\setlength{\tabcolsep}{3.5pt}
\begin{tabular}{lrrrrc}
\toprule
Actor & Cells & $\epsilon=0$ & $1/65536$ & $1/4096$ & Uniform allowance \\
\midrule
CR8 & 64 & 28 & 28 & 28 & -- \\
CB8 & 64 & 64 & 64 & 55 & $1/32768$ \\
CR16 & 256 & 256 & 248 & 195 & $1/262144$ \\
CB16 & 256 & 256 & 243 & 200 & $1/8388608$ \\
\bottomrule\end{tabular}}
\par\smallskip{\footnotesize Counts use the strict inequality D greater than U plus epsilon. U is the exact objective of the attained feasible reference fit. The last column is a dyadic allowance strictly below the smallest positive switch surplus, when such a surplus is certified. A dash is unresolved, not evidence of instability. The independent verifier checks every one of the 20,480 action-switch duals.}
\end{table}

\begin{table}[htbp]
\centering
\caption{Cost comparisons uniform over exact qualification minimizers}\label{tab:optimizer-classes81}
{\small\setlength{\tabcolsep}{3.5pt}
\begin{tabular}{lc}
\toprule
Complete optimizer classes & Original-law cost difference \\
\midrule
CB16 minus CB8 & [-0.003286, -0.002241] \\
CB16 minus CR16 & [-0.002013, -0.000955] \\
CB8 minus CR16 & [0.000748, 0.001812] \\
\bottomrule\end{tabular}}
\par\smallskip{\footnotesize CB8, CR16 and CB16 have fully stable first-date covers at epsilon zero. Their terminal actors are fixed by the common terminal payoff. Consequently each stated exact-minimizer class gives the attained complete policy and inherits its interval without enlargement. This is a deterministic transfer of the same finite statistical family, not new sampling over optimizers.}
\end{table}


At zero extra objective allowance, the two fine-resolution constructions and coarse Bernstein have complete first-date action stability. Every exact minimizer of each of these fixed-label qualification problems therefore gives its attained complete policy, not merely a similar fitted value. The original-law intervals in Table~\ref{tab:optimizer-classes81} extend to these optimizer classes without additional probability expenditure. The same conclusion holds for the positive dyadic sublevels displayed in Table~\ref{tab:stability81}.

The standard $1/65536$ allowance leaves some fine-grid cells unresolved. Theorem~\ref{thm:stability81} then supplies an explicit enlargement based on their uniform initial mass. This conservative enlargement need not preserve a small cross-method sign. Coarse ReLU also retains unresolved cells even at zero allowance. None is relabeled as a discovered instability, and exact-minimizer invariance is not a claim that an arbitrary fitting routine reaches the requisite sublevel. The new statement addresses optimizer-independent policy stability under explicit attained-objective conditions; it does not estimate a hidden-layer success probability.

\subsection{Controller choice and complete costs}
Both cash and conditional marginal and pair intervals can be used together on their original simultaneous event. Exact shortest-path closure gives consistent cost potentials and the pure-choice regret bounds of the inherited decision theorem. The four-term interactions remain separately reported; sharpness of the pair-constraint calculation is not asserted after adding those extra restrictions. The complete nine-controller closure and attaining potentials are deposited in the revision audit.

For the fine calibrated controllers, the positive ReLU-minus-Bernstein interval determines a nonempty region of additional Bernstein charges for which its lower policy cost outweighs its extra charge. It also determines the opposite region in which ReLU is preferred after that charge. All such statements are conditional on an explicitly supplied conversion of computation, installation and deployment into normalized cost units. The exact boundary values may be tightened by the common pair-constraint closure; no new simulation is needed for that deterministic calculation.

The construction work is recorded rather than presumed equal. Legacy clocks include a label-loss ablation that the calibrated services do not execute. Their difference therefore does not identify an end-to-end speed gain from calibration. Exact action and expected-feature evaluation counts are reported alongside the clocks, and the original cost-inference, barrier and publication services retain separate measured costs. The present result supplies an identified policy-cost and conditional purchase comparison in the original economy. It does not extrapolate a two-date fitted-continuation calculation into a measured long-horizon speed result for the recursively queried vector tree.
